\section{引言}\label{sec:intro}
盒马等生鲜电商每天以有限的库存开始营业，营业期间无法补货。顾客打开 App 时，平台需要在几毫秒内决定推荐哪些商品。这一决定会影响后续的销售：一件稀缺商品卖给当前顾客，后面更需要它的顾客可能就买不到了。题目\cite{assignment}给出的算法 1 按单次期望收益 $r_jp_{ij}$ 为每类顾客固定推荐前 $K$ 件商品，没有考虑库存。本文围绕题目的三个问题展开：
\begin{enumerate}[label=(\arabic*)]
\item 有限库存下能否继续使用算法 1？是否存在使其收入损失超过 30\% 的例子？
\item 考虑库存约束时，离线（hindsight）模型是什么？
\item 如何设计满足库存约束、可实时运行的在线推荐算法？
\end{enumerate}

本文的主要工作与结论如下。
\begin{itemize}
\item \textbf{问题一。} 算法 1 会推荐缺货商品，违反库存可行性。本文构造了一族单类型算例，算法 1 相对可行策略的收入损失在销售期 $T\ge14$ 时超过 30\%，$T=20$ 时为 47.08\%，并随 $T$ 增大趋于 100\%。
\item \textbf{问题二。} 建立含库存约束的流体线性规划，证明其最优值是任意在线策略期望收入的上界。课程实例的最优值为 706，商品 1 的影子价格为 9。本文比较了三种离线表达，并给出了库存缩放下影子价格的变化。
\item \textbf{问题三。} 提出 \LPDual{} 策略：每期以当前库存与剩余期数重新求解一次流体线性规划，用库存约束的对偶价格作为机会成本，推荐调整后收益最高的至多 $K$ 件有货商品。策略每期只求解一个 $mn$ 维的连续线性规划，满足零缺货推荐，可以在任意 $m,n,T,K$ 下运行。
\item \textbf{评价。} 本文用库存格点上的前向递推精确计算各策略的期望收入，不含蒙特卡洛误差。五份数据上，\LPDual{} 的平均遗憾为 \GavgDual，最大遗憾为 \GmaxDual，均低于库存感知贪心策略与有限差分版本；遗憾中有 \FluidShare\% 是流体模型与随机系统之间的松弛间隙，与预知全部顾客类型的最优策略相比，策略损失只有 \LavgDual。
\item \textbf{分析。} 通过损失分解、1000 条生成序列、分块到达顺序、库存与推荐位数的灵敏度，以及销售期放大与较大实例，本文说明了各策略在何时、因何损失收入，以及流体近似的适用范围。
\end{itemize}

全文结构如下：第\ref{sec:setting}节描述问题与模型；第\ref{sec:q1}、\ref{sec:q2}、\ref{sec:policy}节依次回答三个问题；第\ref{sec:eval}节给出评价方法；第\ref{sec:exp}节报告数值实验；第\ref{sec:conclusion}节总结。附录说明程序结构与自动化测试。
